\documentclass[10pt, a4paper]{article}

% ---------------------------------------------------------------
% Packages
% ---------------------------------------------------------------
\usepackage[utf8]{inputenc}
\usepackage[T1]{fontenc}
\usepackage[margin=0.80in, top=1.0in, bottom=1.0in]{geometry}
\usepackage{booktabs}
\usepackage{amsmath, amssymb}
\usepackage{xcolor}
\usepackage{hyperref}
\usepackage{fancyhdr}
\usepackage{array}
\usepackage{tcolorbox}
\usepackage{enumitem}
\usepackage{multirow}
\usepackage{makecell}
\usepackage{caption}
\usepackage{tabularx}      % flexible X columns
\usepackage{adjustbox}     % \resizebox / adjustbox env for overflow tables
\usepackage{pdflscape}     % landscape pages for very wide tables

% ---------------------------------------------------------------
% Colours
% ---------------------------------------------------------------
\definecolor{primaryblue}{RGB}{24, 76, 120}
\definecolor{secondaryblue}{RGB}{41, 128, 185}
\definecolor{warnorange}{RGB}{200, 100, 20}
\definecolor{okgreen}{RGB}{30, 130, 60}
\definecolor{lightbg}{RGB}{245, 247, 250}
\definecolor{warnbg}{RGB}{255, 248, 230}

% ---------------------------------------------------------------
% Hyperlinks
% ---------------------------------------------------------------
\hypersetup{colorlinks=true, linkcolor=secondaryblue,
            citecolor=secondaryblue, urlcolor=secondaryblue}

% ---------------------------------------------------------------
% Header / Footer
% ---------------------------------------------------------------
\pagestyle{fancy}
\fancyhf{}
\lhead{\small \textbf{GTE-PPIS: Architecture Audit \& Results Comparison}}
\rhead{\small \textbf{Report Date: September 3, 2026}}
\lfoot{\small Logs: \texttt{noesmtest\_25thAugust\_13.log} \quad \texttt{rsa\_features\_test.log}}
\rfoot{\small Page \thepage}
\renewcommand{\headrulewidth}{0.5pt}
\renewcommand{\footrulewidth}{0.5pt}

% ---------------------------------------------------------------
% Formatting
% ---------------------------------------------------------------
\setlength{\parindent}{0pt}
\setlength{\parskip}{5pt}
\captionsetup{font=small, labelfont=bf}

% tabularx: make X columns use \footnotesize automatically
\newcolumntype{Y}{>{\footnotesize\raggedright\arraybackslash}X}

% ---------------------------------------------------------------
% Helper commands
% ---------------------------------------------------------------
\newcommand{\match}[1]{\textcolor{okgreen}{\textbf{#1}}}
\newcommand{\mismatch}[1]{\textcolor{warnorange}{\textbf{#1}}}
\newcommand{\warn}[1]{\textcolor{warnorange}{\textbf{#1}}}
\newcommand{\code}[1]{\texttt{\small #1}}

% ================================================================
\begin{document}
% ================================================================

% ---------------------------------------------------------------
% Title
% ---------------------------------------------------------------
\begin{center}
    {\LARGE \bfseries \color{primaryblue} GTE-PPIS: Architecture Audit \& Results Comparison}\\[4pt]
    {\large \color{secondaryblue}
        Codebase vs.\ Published Paper $\cdot$
        Reproduction Gap Analysis $\cdot$
        Three-Way Results Table}\\[6pt]
    \rule{\textwidth}{1pt}\\[5pt]
    \textbf{Logs used:}
    \code{noesmtest\_25thAugust\_13.log} (no-ESM2 baseline, Aug 25 2026) \quad|\quad
    \code{rsa\_features\_test.log} (gated ESM-2, Aug 13 2026)
    \rule{\textwidth}{0.5pt}
\end{center}

\vspace{4pt}

% ================================================================
\section{Architecture Comparison: Paper vs.\ Codebase}
% ================================================================

The following subsections compare each architectural component against
what is described in the published GTE-PPIS paper.
\match{Green} = confirmed match; \mismatch{Orange} = divergence or addition.

% ---------------------------------------------------------------
\subsection{EGNN Branch}
% ---------------------------------------------------------------

% Four-column wide table: use tabularx with three Y (flexible) columns + fixed Status column
\begin{table}[ht]
\caption{EGNN branch: paper description vs.\ codebase (\code{EGNN\_model.py}, \code{final\_model.py})}
\label{tab:egnn_compare}
\begin{tabularx}{\linewidth}{>{\footnotesize\bfseries\raggedright\arraybackslash}p{2.2cm} Y Y
    >{\footnotesize\centering\arraybackslash}p{2.8cm}}
\toprule
\normalsize\textbf{Dimension} & \normalsize\textbf{Paper} & \normalsize\textbf{Codebase} & \normalsize\textbf{Status} \\
\midrule
Layer count &
  Not explicitly stated; ablation ``NoEGNN'' implies integral EGNN &
  \code{n\_layers=10} (hardcoded in \code{FinalModel}) &
  \warn{Unconfirmed} \\
\addlinespace
Hidden dim &
  Not reported; implicit from overall model width &
  \code{hidden\_nf = HIDDEN\_DIM = 256} &
  \warn{Unconfirmed} \\
\addlinespace
Residual &
  ``NoResidual'' ablation (Table~4): AUROC 0.871$\to$0.852, implying residuals are active &
  \code{residual=True} in every \code{E\_GCL}; skip: $h_i^{l+1} = h_i^l + \phi_h(\cdot)$ &
  \match{Match in intent; previously False --- now fixed} \\
\addlinespace
Normalization &
  Not described for EGNN; norm discussion is specific to GT branch &
  \textbf{No LayerNorm or BatchNorm inside \code{E\_GCL}}; only MLP activations (SiLU) &
  \mismatch{Paper ambiguous; no inter-layer norm} \\
\addlinespace
Coordinate update &
  Standard EGNN equivariant update &
  \code{coords\_agg='mean'}, \code{tanh=True} (bounds update; our stability addition) &
  \mismatch{tanh=True is our addition} \\
\addlinespace
Attention gate &
  Not reported &
  \code{attention=True} (per-edge sigmoid gate) &
  \mismatch{Our addition; paper omits} \\
\addlinespace
Edge features &
  Distance + cosine (2-dim) &
  \code{in\_edge\_nf=2}; same 2-dim &
  \match{Match} \\
\bottomrule
\end{tabularx}
\end{table}

\begin{tcolorbox}[colback=warnbg, colframe=warnorange, title=\warn{EGNN Key Concern}]
The paper does not specify EGNN layer count or hidden dim numerically.
Our codebase uses \textbf{10 EGNN layers with hidden dim 256}.
Additionally, \code{tanh=True} and \code{attention=True} are our own additions not present in the
original, which could affect gradient flow and partially explain the reproduction gap.
The \code{residual=True} fix aligns with the paper's ablation evidence (Table~4).
\end{tcolorbox}

% ---------------------------------------------------------------
\subsection{Graph Transformer Branch}
% ---------------------------------------------------------------

\begin{table}[ht]
\caption{Graph Transformer branch: paper Eqs.\ 7--13 vs.\ \code{GraphTransformer\_Block.py}}
\label{tab:gt_compare}
\begin{tabularx}{\linewidth}{>{\footnotesize\bfseries\raggedright\arraybackslash}p{2.6cm} Y Y
    >{\footnotesize\centering\arraybackslash}p{2.8cm}}
\toprule
\normalsize\textbf{Dimension} & \normalsize\textbf{Paper (Eqs.\ 7--13)} & \normalsize\textbf{Codebase} & \normalsize\textbf{Status} \\
\midrule
Layer count &
  4 layers cited &
  \code{num\_layers=4} (3$\times$\code{GraphTransformerModule} + 1$\times$\code{FinalGraphTransformerModule}) &
  \match{Match} \\
\addlinespace
Attention heads & 4 heads & \code{num\_attention\_heads=4} & \match{Match} \\
\addlinespace
Normalization &
  ``layer normalization'' &
  \code{norm\_to\_apply='layer'} $\Rightarrow$ \code{nn.LayerNorm} &
  \match{Match} \\
\addlinespace
LayerNorm placement &
  Pre-norm before MHA (Eq.\ 7), post-residual before FFN (Eq.\ 9) &
  LN1 before MHA; LN2 after first residual, before FFN MLP --- exact pre-norm pattern &
  \match{Match} \\
\addlinespace
GT residual (module) &
  Eq.\ 8: $h_i^{(l+1)} = h_i^{(l)} + O(\hat{h}_i)$; two residuals per layer &
  Two residual connections in \code{run\_gt\_layer} (\code{node\_feats\_in1}, \code{node\_feats\_in2}) &
  \match{Match in module} \\
\addlinespace
Edge feature update &
  Intermediate layers update edge reps; final layer does not &
  \code{update\_edge\_feats=True} for intermediate, \code{False} for \code{FinalGTModule} &
  \match{Match} \\
\addlinespace
\textbf{GT residual at call site} &
  Residuals enabled &
  \code{transformer\_residual=False} passed from \code{FinalModel} line~56 &
  \mismatch{\textbf{DIVERGENCE}: residual \emph{disabled} in all GT layers} \\
\addlinespace
Hidden dim &
  Not stated &
  \code{num\_hidden\_channels=128} &
  \warn{Unconfirmed; paper may use 256} \\
\addlinespace
Output MLP &
  2-class linear head (Eq.\ 13) &
  \code{output\_linner = nn.Linear(128, 2)} &
  \match{Match} \\
\bottomrule
\end{tabularx}
\end{table}

\begin{tcolorbox}[colback=warnbg, colframe=warnorange,
    title=\warn{Critical GT Divergence: \code{transformer\_residual=False}}]
In \code{final\_model.py} line~56:
\begin{verbatim}
self.GT = GraghTransformer(..., transformer_residual=False)
\end{verbatim}
This disables \emph{both} residual connections inside every GT layer.
The paper explicitly uses residual connections (Eqs.\ 8 and 10), and the ablation in Table~4
shows ``NoResidual'' reduces CV AUROC by 1.9 points (0.871 $\to$ 0.852).
\textbf{This is the highest-impact confirmed divergence and the most likely
primary contributor to the reproduction gap.}
\end{tcolorbox}

% ---------------------------------------------------------------
\subsection{Feature Set}
% ---------------------------------------------------------------

\begin{table}[ht]
\caption{Node and edge feature comparison: paper vs.\ \code{data\_generator.py}}
\label{tab:features}
\begin{tabularx}{\linewidth}{>{\footnotesize\bfseries\raggedright\arraybackslash}p{2.2cm}
    Y Y
    >{\footnotesize\centering\arraybackslash}p{0.8cm}
    >{\footnotesize\centering\arraybackslash}p{1.5cm}}
\toprule
\normalsize\textbf{Feature} & \normalsize\textbf{Paper} & \normalsize\textbf{Codebase} & \normalsize\textbf{Dim} & \normalsize\textbf{Status} \\
\midrule
PSSM          & 20-dim per-residue PSSM profile & \code{Feature/pssm/*.npy} & 20 & \match{Match} \\
HMM           & 20-dim HMM profile              & \code{Feature/hmm/*.npy}  & 20 & \match{Match} \\
DSSP          & 14-dim secondary struct / solvent accessibility & \code{Feature/dssp/*.npy} & 14 & \match{Match} \\
resAF         & 7-dim residue atom features     & \code{Feature/resAF/*.npy} & 7  & \match{Match} \\
Total node    & 61-dim                          & \code{INPUT\_DIM = 61}    & 61 & \match{Match} \\
\addlinespace
Edge: distance & Normalized Euclidean, 14\,\AA\ cutoff & \code{MAP\_CUTOFF=14}, normalized by \code{DIST\_NORM=15} & 1 & \match{Match} \\
Edge: cosine  & Cosine similarity of $C_\alpha$ pseudo-positions & \code{cossim(...)} rescaled to $[0,1]$ & 1 & \match{Match} \\
Total edge    & 2 & \code{in\_edge\_nf=2} & 2 & \match{Match} \\
\midrule
\multicolumn{5}{l}{\footnotesize\textit{Our Extensions (not in paper):}} \\
\addlinespace
ESM-2 embeddings & --- & \code{Feature/esm2/*.npy}, ESM-2-650M, fp16 cached & 1280 & \mismatch{Added} \\
RSA features  & RSA embedded inside DSSP 14-dim block & Separate \code{Feature/rsa/*.npy} continuous scalar; gate supervision & 1 & \mismatch{Added} \\
Fusion module & --- & \code{FeatureFusionModule}: gated / concat / cross-attn & --- & \mismatch{Added} \\
\bottomrule
\end{tabularx}
\end{table}

The paper's train set is \textbf{Train\_334} (334 chains, 10,336 interacting residues, 15.61\% positive).
Our codebase loads \code{Train\_335.pkl} and pops \code{`2j3rA'}, giving 334 effective chains.
The measured positive ratio of 15.61\% confirms the feature set is correctly reproduced.

% ---------------------------------------------------------------
\subsection{Loss Function}
% ---------------------------------------------------------------

\begin{table}[ht]
\caption{Loss function: paper vs.\ \code{loss.py} + \code{final\_model.py}}
\label{tab:loss}
\begin{tabularx}{\linewidth}{>{\footnotesize\bfseries\raggedright\arraybackslash}p{2.5cm}
    Y Y
    >{\footnotesize\centering\arraybackslash}p{2.4cm}}
\toprule
\normalsize\textbf{Aspect} & \normalsize\textbf{Paper} & \normalsize\textbf{Codebase} & \normalsize\textbf{Status} \\
\midrule
Loss type &
  Cross-entropy (no class weighting or focal mentioned) &
  \code{FocalLoss(alpha=[1.0, neg/pos], gamma=2.0)} &
  \mismatch{\textbf{Divergence}} \\
Class weighting &
  Not mentioned &
  \code{alpha=[1.0, 5.4056]} (neg/pos ratio from Train\_335) &
  \mismatch{Added} \\
Focal modulation &
  Not used &
  $\gamma=2.0$ (down-weights easy examples) &
  \mismatch{Added} \\
Auxiliary losses &
  None &
  \code{loss\_gate} (RSA MSE, $\lambda=0.1$) + \code{loss\_agree} (branch MSE, $\lambda=0.1$) &
  \mismatch{Added (gated only)} \\
Gradient clipping &
  Not mentioned &
  \code{clip\_grad\_norm\_(max\_norm=1.0)} &
  \mismatch{Added} \\
\bottomrule
\end{tabularx}
\end{table}

\begin{tcolorbox}[colback=warnbg, colframe=warnorange, title=\warn{Loss Divergence Note}]
The paper uses standard cross-entropy.
Our no-ESM2 run already uses \code{FocalLoss($\gamma$=2.0)} with class weighting.
Focal Loss preferentially trains on hard examples and changes the effective gradient landscape ---
a \textbf{second confirmed divergence} that directly affects the learned decision boundary.
\end{tcolorbox}

% ---------------------------------------------------------------
\subsection{Training Setup}
% ---------------------------------------------------------------

\begin{table}[ht]
\caption{Training hyper-parameters: paper vs.\ \code{train.py} + \code{data\_generator.py}}
\label{tab:training}
\begin{tabularx}{\linewidth}{>{\footnotesize\bfseries\raggedright\arraybackslash}p{2.6cm}
    Y Y
    >{\footnotesize\centering\arraybackslash}p{2.6cm}}
\toprule
\normalsize\textbf{Hyper-parameter} & \normalsize\textbf{Paper} & \normalsize\textbf{Codebase} & \normalsize\textbf{Status} \\
\midrule
Batch size   & 1 (per-protein)   & \code{BATCH\_SIZE = 1}   & \match{Match} \\
Epochs       & 50                & \code{NUMBER\_EPOCHS = 50} & \match{Match} \\
Optimizer    & Adam              & \code{torch.optim.Adam}  & \match{Match} \\
Learning rate & $10^{-3}$        & \code{LEARNING\_RATE = 1e-3} & \match{Match} \\
Weight decay & 0                 & \code{WEIGHT\_DECAY = 0} & \match{Match} \\
LR schedule  & Not described &
  \code{ReduceLROnPlateau(mode='max', factor=0.6, patience=5)} on valid AUPRC &
  \mismatch{Addition; not in paper} \\
Dropout      & Not stated; GT uses dropout & \code{DROPOUT = 0.1} inside GT FFN MLP & \match{Consistent} \\
CV strategy  & 5-fold on Train\_334 &
  \code{KFold(n\_splits=5, shuffle=True)} with \code{SEED=2024} &
  \match{Match (seed may differ)} \\
Random seed  & Not stated        & \code{SEED = 2024} & \warn{Unknown if paper seed differs} \\
Checkpoint criterion & Best by validation AUPRC &
  \code{best\_val\_aupr < result\_valid['AUPRC']} &
  \match{Match} \\
\textbf{Threshold scan} &
  Includes threshold=0.0 (causes MCC=0 collapse) &
  Starts from 0.01 (\code{range(1,100)/100}); excludes degenerate all-positive case &
  \mismatch{Bug-fix divergence} \\
\bottomrule
\end{tabularx}
\end{table}

\newpage

% ================================================================
\section{Results Comparison}
% ================================================================

\subsection{Reproduced Paper Tables (as Published)}

\begin{table}[ht]
\caption{Paper Table~2 extended --- Evaluation on \texttt{Test\_60} (published numbers + our runs).
$\ddagger$ = single run, no multi-seed verification.}
\label{tab:paper_test60}
\resizebox{\linewidth}{!}{%
\begin{tabular}{lrrrrrrr}
\toprule
\textbf{Method} & \textbf{ACC} & \textbf{Prec.} & \textbf{Recall} & \textbf{F1} & \textbf{AUROC} & \textbf{MCC} & \textbf{AUPRC} \\
\midrule
ScanNet    & 0.681 & 0.245 & 0.547 & 0.339 & 0.684 & 0.191 & 0.282 \\
DELPHI     & 0.687 & 0.270 & 0.577 & 0.368 & 0.691 & 0.220 & 0.307 \\
HN-PPISP   & 0.738 & 0.278 & 0.415 & 0.333 & 0.656 & 0.184 & 0.281 \\
DeepPPISP  & 0.750 & 0.297 & 0.428 & 0.351 & 0.676 & 0.207 & 0.295 \\
MaSIF-site & 0.780 & 0.370 & 0.561 & 0.446 & 0.775 & 0.379 & 0.372 \\
EDLMPPI    & 0.781 & 0.360 & 0.503 & 0.420 & 0.748 & 0.295 & 0.379 \\
GraphPPIS  & 0.792 & 0.389 & 0.558 & 0.458 & 0.784 & 0.343 & 0.430 \\
RGCNPPIS   & 0.799 & 0.404 & 0.572 & 0.474 & 0.803 & 0.439 & 0.471 \\
AGAT-PPIS  & 0.856 & 0.539 & 0.603 & 0.569 & 0.867 & 0.484 & 0.574 \\
\midrule
\textbf{GTE-PPIS (paper)}           & \textbf{0.861} & \textbf{0.557} & \textbf{0.611} & \textbf{0.582} & \textbf{0.873} & \textbf{0.500} & \textbf{0.611} \\
\midrule
\textit{Ours (A): No-ESM2}          & 0.8405 & 0.4978 & 0.5734 & 0.5319 & 0.8476 & 0.4388 & 0.5380 \\
\textit{Ours (B): Gated ESM-2$^{\ddagger}$} & 0.8356 & 0.4841 & 0.5395 & 0.5089 & 0.8236 & 0.4125 & 0.5056 \\
\bottomrule
\end{tabular}}
\end{table}

\begin{table}[ht]
\caption{Paper Table~3 extended --- Evaluation on \texttt{Test\_315-28}, \texttt{Btest\_25}, and \texttt{UBtest\_25}/\texttt{31-6} (published numbers + our runs).
$\ddagger$ = single run, no multi-seed verification.
$\dagger$ = our split (\texttt{UBtest\_31-6}) differs from paper's \texttt{UBtest\_25}; not directly comparable.}
\label{tab:paper_test315}
\resizebox{\linewidth}{!}{%
\begin{tabular}{lcccccc}
\toprule
\multirow{2}{*}{\textbf{Method}} &
\multicolumn{2}{c}{\textbf{Test\_315-28}} &
\multicolumn{2}{c}{\textbf{Btest\_25}} &
\multicolumn{2}{c}{\textbf{UBtest\_25 / 31-6}} \\
\cmidrule(lr){2-3}\cmidrule(lr){4-5}\cmidrule(lr){6-7}
& \textbf{MCC} & \textbf{AUPRC} & \textbf{MCC} & \textbf{AUPRC} & \textbf{MCC} & \textbf{AUPRC} \\
\midrule
GraphPPIS  & 0.335 & 0.408 & 0.339 & 0.381 & 0.298 & 0.330 \\
RGCNPPIS   & 0.352 & 0.420 & 0.375 & 0.400 & 0.296 & \textbf{0.354} \\
AGAT-PPIS  & 0.442 & 0.525 & 0.440 & 0.511 & 0.301 & 0.325 \\
\midrule
\textbf{GTE-PPIS (paper)}           & \textbf{0.511} & \textbf{0.598} & \textbf{0.471} & \textbf{0.545} & \textbf{0.320} & 0.343 \\
\midrule
\textit{Ours (A): No-ESM2}          & 0.4354 & 0.5182 & --- & --- & 0.2713$^{\dagger}$ & 0.3115$^{\dagger}$ \\
\textit{Ours (B): Gated ESM-2$^{\ddagger}$} & 0.3601 & 0.4381 & --- & --- & 0.3706$^{\dagger}$ & 0.4182$^{\dagger}$ \\
\bottomrule
\end{tabular}}
\end{table}

\textit{Note: the paper evaluates on \textbf{UBtest\_25} (25 proteins, 711 interacting residues).
Our codebase uses \textbf{UBtest\_31-6}, which is \textbf{not directly comparable}.}

% ---------------------------------------------------------------
\subsection{Three-Way Results Comparison}
% ---------------------------------------------------------------

% 9 columns: use \resizebox to guarantee it fits
\begin{table}[ht]
\caption{Three-way results comparison (5-fold CV averages).
Bold = best among our experimental runs.
$\ddagger$ = single run, no multi-seed verification.
$\dagger$ = different test split from paper.}
\label{tab:three_way}
\resizebox{\linewidth}{!}{%
\begin{tabular}{llccccccc}
\toprule
\textbf{Test Set} & \textbf{System} & \textbf{AUROC} & \textbf{AUPRC} & \textbf{F1} & \textbf{MCC} & \textbf{Precision} & \textbf{Recall} & \textbf{ACC} \\
\midrule
\multirow{3}{*}{\texttt{Test\_60}}
& (A) No-ESM2 (ours)         & \textbf{0.8476} & \textbf{0.5380} & \textbf{0.5319} & \textbf{0.4388} & \textbf{0.4978} & \textbf{0.5734} & \textbf{0.8405} \\
& (B) Gated ESM-2$^\ddagger$ & 0.8236 & 0.5056 & 0.5089 & 0.4125 & 0.4841 & 0.5395 & 0.8356 \\
& (C) Paper (GTE-PPIS)       & 0.873  & 0.611  & 0.582  & 0.500  & 0.557  & 0.611  & 0.861  \\
\midrule
\multirow{3}{*}{\texttt{Test\_315-28}}
& (A) No-ESM2 (ours)         & \textbf{0.8585} & \textbf{0.5182} & \textbf{0.5189} & \textbf{0.4354} & 0.4181 & \textbf{0.6847} & \textbf{0.8199} \\
& (B) Gated ESM-2$^\ddagger$ & 0.8117 & 0.4381 & 0.4597 & 0.3601 & 0.3727 & 0.6012 & 0.7994 \\
& (C) Paper (GTE-PPIS)       & ---    & 0.598  & ---    & 0.511  & ---    & ---    & ---    \\
\midrule
\multirow{3}{*}{\texttt{UBtest\_31-6}$^\dagger$}
& (A) No-ESM2 (ours)         & 0.7611 & 0.3115 & 0.3542 & 0.2713 & 0.3717 & 0.3418 & 0.8504 \\
& (B) Gated ESM-2$^\ddagger$ & \textbf{0.8191} & \textbf{0.4182} & \textbf{0.4454} & \textbf{0.3706} & 0.4418 & 0.4554 & \textbf{0.8646} \\
& (C) Paper (\texttt{UBtest\_25}) & --- & 0.343 & --- & 0.320 & ---    & ---    & ---    \\
\bottomrule
\end{tabular}}
\end{table}

% ---------------------------------------------------------------
\subsection{Reproduction Gap: No-ESM2 (A) vs.\ Paper (C) on \texttt{Test\_60}}
% ---------------------------------------------------------------

\begin{table}[ht]
\caption{Absolute and relative metric gaps on \texttt{Test\_60}}
\label{tab:gap}
\resizebox{\linewidth}{!}{%
\begin{tabular}{lccccccc}
\toprule
& \textbf{AUROC} & \textbf{AUPRC} & \textbf{F1} & \textbf{MCC} & \textbf{Precision} & \textbf{Recall} & \textbf{ACC} \\
\midrule
Paper (C)    & 0.8730 & 0.6110 & 0.5820 & 0.5000 & 0.5570 & 0.6110 & 0.8610 \\
Ours (A)     & 0.8476 & 0.5380 & 0.5319 & 0.4388 & 0.4978 & 0.5734 & 0.8405 \\
\midrule
$\Delta$ (abs) &
  \mismatch{$-$0.0254} & \mismatch{$-$0.0730} & \mismatch{$-$0.0501} &
  \mismatch{$-$0.0612} & \mismatch{$-$0.0592} & \mismatch{$-$0.0376} & \mismatch{$-$0.0205} \\
Rel.\ gap (\%) &
  \mismatch{$-$2.9} & \mismatch{$-$12.0} & \mismatch{$-$8.6} &
  \mismatch{$-$12.2} & \mismatch{$-$10.6} & \mismatch{$-$6.2} & \mismatch{$-$2.4} \\
\bottomrule
\end{tabular}}
\end{table}

The \textbf{MCC gap of $-0.0612$ (12.2\% relative)} and
\textbf{AUPRC gap of $-0.073$ (12.0\% relative)} cannot be attributed to seed variance
alone ($<$1--2\% in typical reproducibility studies on this scale).

% ================================================================
\section{Divergence Ranking by Estimated Impact on Reproduction Gap}
% ================================================================

\begin{table}[ht]
\caption{Ranked divergences from the published GTE-PPIS design, ordered by estimated contribution to the reproduction gap}
\label{tab:divergences}
\begin{tabularx}{\linewidth}{
    >{\footnotesize\centering\arraybackslash}p{0.9cm}
    >{\footnotesize\bfseries\raggedright\arraybackslash}p{2.2cm}
    Y
    >{\footnotesize\raggedright\arraybackslash}p{3.5cm}}
\toprule
\normalsize\textbf{Rank} & \normalsize\textbf{Component} &
\normalsize\textbf{Divergence Description} & \normalsize\textbf{Estimated Impact} \\
\midrule
\textbf{1} & GT Branch &
  \code{transformer\_residual=False} disables both skip connections in all 4 GT layers;
  paper Table~4 ``NoResidual'' shows whole-model $\Delta$AUROC = $-$12.9\% on Test\_60 &
  \textbf{High (weak evidence only)} --- residuals matter somewhere in the model is confirmed;
  attribution to GT branch specifically is an \emph{untested inference} (see Section~4) \\
\addlinespace
\textbf{2} & Loss function &
  Focal Loss ($\gamma$=2.0) + class-weight alpha replacing plain CE;
  changes gradient weighting across the entire training distribution &
  \textbf{High} --- shifts precision/recall trade-off; precision gap is 5.9 pts \\
\addlinespace
\textbf{3} & EGNN modifications &
  \code{tanh=True} bounds coordinate updates; \code{attention=True} adds per-edge gating
  (both absent in paper) &
  \textbf{Medium} --- stabilizes training but restricts model expressivity \\
\addlinespace
\textbf{4} & LR schedule &
  \code{ReduceLROnPlateau} on validation AUPRC; paper does not describe any schedule &
  \textbf{Low--Medium} --- may cause early LR reductions in some folds \\
\addlinespace
\textbf{5} & Random seed &
  Paper seed unknown; ours is 2024; 5-fold splits will differ &
  \textbf{Low} --- $<$2\% variation typical \\
\addlinespace
\textbf{6} & Threshold scan fix &
  Start from 0.01 instead of 0.0; eliminates degenerate recall=1 solutions &
  \textbf{Positive correction} --- this is a bug-fix; \emph{reduces} metric inflation,
  not a source of the gap \\
\addlinespace
\textbf{7} & EGNN layer count &
  10 layers vs.\ unspecified; if paper uses fewer, deeper network may over-regularize &
  \textbf{Unknown} --- cannot quantify without paper's released code \\
\bottomrule
\end{tabularx}
\end{table}

\begin{tcolorbox}[colback=lightbg, colframe=primaryblue,
    title=\textbf{Root Cause Hypothesis and Recommended Action}]
The most actionable hypothesis is that the \textbf{12\% relative MCC gap} is
predominantly explained by \textbf{two confirmed divergences}:

\begin{enumerate}[leftmargin=1.5em, itemsep=2pt]
    \item \textbf{(Primary)} \code{transformer\_residual=False} in our GT call ---
    disabling residuals directly weakens representational capacity, consistent with
    the paper's own ablation evidence. \textbf{Fix}: change to \code{transformer\_residual=True}.

    \item \textbf{(Secondary)} Focal Loss with class-weight alpha replacing plain CE ---
    originally introduced to compensate for EGNN training collapse (itself caused by
    the now-fixed \code{residual=False} in E\_GCL).
    Continued use shifts the effective decision boundary vs.\ the paper.
    \textbf{Fix}: retrain with \code{focal\_gamma=0.0} (reduces to weighted CE) or plain CE
    once the GT residual is restored.
\end{enumerate}

\textbf{Recommended next experiment}: retrain with \code{transformer\_residual=True}
and \code{focal\_gamma=0.0} to isolate each contribution.
Do \textbf{not} claim any single-run result from set~(B) as an improvement without
multi-seed replication ($\geq$3 seeds).
\end{tcolorbox}


\newpage

% ================================================================
\section{Addendum: Verification of Two Key Claims}
% ================================================================

% ---------------------------------------------------------------
\subsection{Finding 1 --- Scope of the ``NoResidual'' ablation (paper Table~4)}
% ---------------------------------------------------------------

\paragraph{What the paper actually says.}
Section ``Model architecture analysis'' (PMC12199915, paragraph starting at Table~4)
defines the ablation variant as:

\begin{quote}
\small\itshape
``\ldots\, NoResidual (\textbf{excluding the residual connections}) \ldots''
\end{quote}

\noindent
The definition uses the unqualified plural. The paper gives no further specification of
\emph{which branch's} residuals were removed. There is no branch-isolated ablation variant.

\paragraph{Exact numbers from Table~4.}
NoResidual causes: CV: $\Delta$AUROC $= -2.2\%$, $\Delta$AUPRC $= -8.4\%$;
Test\_60: $\Delta$AUROC $= -12.9\%$, $\Delta$AUPRC $= -36.5\%$.
These are measured for the whole model (both branches with residuals removed simultaneously).

\paragraph{What the prior report established --- and what it did not.}
The report used these numbers as evidence that \code{transformer\_residual=False} in
the GT call site accounts for the 12\% MCC gap. That reasoning conflates two separate claims:

\begin{enumerate}[noitemsep, leftmargin=1.5em]
    \item \textbf{Confirmed (weak claim)}: Residual connections matter somewhere
    in the full GTE-PPIS model. The NoResidual ablation confirms this clearly.

    \item \textbf{Not confirmed (strong claim)}: That our \code{transformer\_residual=False}
    at only the GT branch call site in \code{final\_model.py} (line~56) is the dominant
    cause of the 12\% MCC gap. The paper provides zero branch-isolated measurement.
    This is an inference, not a measurement.
\end{enumerate}

\begin{tcolorbox}[colback=warnbg, colframe=warnorange,
    title=\warn{Status of Rank 1 Divergence --- Awaiting Retraining}]
\code{transformer\_residual=False} is a genuine divergence from the paper's design
and very likely contributes to the reproduction gap. However, \textbf{the magnitude
of that contribution is unknown} until a retrain with \code{transformer\_residual=True}
is completed. The paper's NoResidual ablation cannot be used to isolate it because
the ablation removes residuals from both branches simultaneously. The strong causal
attribution in the prior version of the report must be treated as an hypothesis until
retraining confirms or refutes it.
\end{tcolorbox}

% ---------------------------------------------------------------
\subsection{Finding 2 --- RSA feature redundancy with DSSP dimension~4}
% ---------------------------------------------------------------

\paragraph{Paper's DSSP feature description.}
Section ``Protein representation / DSSP'' (paper line~274) specifies the 14-dimensional
DSSP block as:

\begin{center}
\small
\begin{tabular}{rll}
\toprule
\textbf{Dims (0-indexed)} & \textbf{Content} & \textbf{Nature} \\
\midrule
0--8   & Secondary structure one-hot (9 DSSP states)         & Discrete   \\
9--12  & $\sin\phi,\cos\phi,\sin\psi,\cos\psi$ (dihedral angles) & Continuous \\
13     & ASA normalized by residue type $\to$ RSA             & Continuous \\
\bottomrule
\end{tabular}
\end{center}

\noindent
Per the paper: ``absolute solvent accessible surface area values were normalized by
residue type to yield relative solvent accessibility (one dimension)'' [paper line~274].

\paragraph{Empirical identification of the RSA column in our DSSP arrays.}
Loading \code{Feature/dssp/1acbI.npy} (shape $63\times14$) shows that \textbf{column~4},
not column~13, carries the RSA signal in our saved files. Columns 9 and 13 have zero
variance for this protein (constant zero), while column~4 has values in $[0,1]$.
Computing Pearson $r$ between each DSSP column and our \code{Feature/rsa/*.npy}
(freesasa-derived) values over protein 1acbI gives $r_{col4} = 0.9858$, $p < 10^{-48}$.

\paragraph{Cross-protein confirmation.}
Computed over 100 proteins with both files present:

\begin{center}
\small
\begin{tabular}{lrrrr}
\toprule
\textbf{Comparison} & \textbf{Mean $r$} & \textbf{Median $r$} & \textbf{Min $r$} & \textbf{Max $r$} \\
\midrule
DSSP col~4 vs.\ \code{Feature/rsa/*.npy} (freesasa) & 0.985 & 0.987 & 0.945 & 0.991 \\
\bottomrule
\end{tabular}
\end{center}

\paragraph{How the two RSA computations differ.}
\begin{itemize}[noitemsep, leftmargin=1.5em]
    \item \textbf{DSSP RSA (col~4 of \code{Feature/dssp/*.npy})}: computed by the DSSP
    algorithm's internal SASA routine, normalized by its own per-residue reference values.
    This is concatenated into the 61-dim node feature vector at
    \code{data\_generator.py} lines~51 and~64, and thus fed as a \emph{model input}
    to both EGNN and GT branches during training and inference.

    \item \textbf{freesasa RSA (\code{Feature/rsa/*.npy})}: computed by
    \code{generate\_rsa\_features.py} using \code{freesasa.calc()} with the
    Shrake--Rupley algorithm and the Miller et al.\ reference $A_{\max}$ table
    (lines~7--12, 53--65 of that script). This value is \textbf{not} concatenated
    into the node feature vector. It is passed only as \code{rsa\_features}
    to \code{compute\_auxiliary\_losses()} (\code{train.py} lines~86--90) as the
    supervision target for the biophysics-gated regularization loss
    $\mathcal{L}_{\text{gate}} = \lambda_{\text{gate}}\cdot\text{MSE}(g_i,\,\text{RSA}_i)$.
\end{itemize}

\begin{tcolorbox}[colback=lightbg, colframe=primaryblue,
    title=\textbf{Redundancy Assessment and Implication for Gate Interpretability}]
The two RSA features measure the same physical quantity ($r = 0.985$ cross-protein)
and differ only in tool (DSSP vs.\ freesasa) and reference $A_{\max}$ table.

\medskip
\textbf{No input redundancy exists}: the freesasa RSA is never added to the input
feature space, so it does not duplicate an input channel.

\medskip
\textbf{Gate supervision is partially redundant}: the gate regularization loss trains
$g_i$ (the gated fusion weight) to track a quantity ($r = 0.985$) that already exists
inside the node features the model consumes. The auxiliary loss therefore provides
much less independent gradient signal than if the RSA target were computed from a
genuinely orthogonal source (e.g., a different structural conformation or a predicted RSA
from sequence alone). The gate likely learns the RSA correlation largely via the input
features, and the auxiliary loss reinforces but does not introduce that correlation.

\medskip
\textbf{Consequence for gate interpretability}: the observation that ``the gate tracks
solvent exposure'' remains valid as a description of model behaviour, but it cannot be
attributed uniquely to the freesasa-based supervision signal --- the gate would likely
learn this correlation anyway from DSSP col~4. The supervision loss contribution is
real but modest.
\end{tcolorbox}

\end{document}
